%! Rates of Change in Polar Functions — Unit 3, Lesson 3.15 — Lesson Plan
\documentclass[10pt]{article}
\usepackage{precalculus-article}
\usepackage{precalculus-boxes}
\usepackage{graphicx}
\graphicspath{{images/}}
\newcommand{\TallMath}[1]{$\displaystyle #1\rule[-1.4em]{0pt}{3.2em}$}
\newcommand{\CourseName}{Precalculus}
\newcommand{\SchoolYear}{2026--2027}
\newcommand{\MeetingLength}{55 minutes}
\newcommand{\UnitNumberName}{Unit 3: Trigonometric and Polar Functions \quad}
\newcommand{\LessonNumberName}{Lesson 3.15: Rates of Change in Polar Functions}

\begin{document}

%----------------------------------------------------------------------
% Title Block
%----------------------------------------------------------------------
\begin{center}
  {\Large\bfseries \CourseName: \SchoolYear} \\[4pt]
  {\small\bfseries \UnitNumberName \LessonNumberName \quad (2--3 instructional periods, \MeetingLength\ each)}
\end{center}

%----------------------------------------------------------------------
% Primary Objective
%----------------------------------------------------------------------
\begin{tcolorbox}[colback=sky, colframe=navy, boxrule=0.9pt, arc=2mm,
    left=2mm, right=2mm, top=1.5mm, bottom=1.5mm]
  \textbf{Primary Objective:} Students will describe how the radius $r$ of a polar
  function changes as the angle $\theta$ increases by computing average rates of change
  of $r$ with respect to $\theta$ and by classifying intervals according to the sign
  and direction of the function. \textit{(Big Idea: Change)}
\end{tcolorbox}

\vspace{0.08in}

%----------------------------------------------------------------------
% Priority Ideas & Skills
%----------------------------------------------------------------------
\begin{skillbox}[Priority Ideas \& Skills]{goldbox}
\begin{tabularx}{\linewidth}{@{} p{0.46\linewidth} X @{}}
\textbf{AP Skills} & \textbf{Key Understandings (EKs)} \\
\midrule
\textbf{2.B} \textit{Construct equivalent representations} ---
  Move between a table of $(θ, r)$ values and a description of whether
  the curve is moving toward or away from the origin.
&
\textbf{EK 3.15.A.1}: The average rate of change of a polar function
$r = f(\theta)$ over $[\theta_1, \theta_2]$ is
$\dfrac{f(\theta_2) - f(\theta_1)}{\theta_2 - \theta_1}$, representing
how quickly $r$ changes per radian. \\[6pt]

\textbf{3.A} \textit{Describe characteristics} ---
  Identify intervals on which $r$ is positive/negative and
  increasing/decreasing; connect these to curve behavior in the plane.
&
\textbf{EK 3.15.A.2}: On intervals where $r > 0$ and increasing, the
curve moves away from the origin; where $r > 0$ and decreasing, it
moves toward the origin. Where $r < 0$, the curve is traced in the
direction opposite the terminal ray. \\[6pt]

&
\textbf{EK 3.15.A.3}: The combination of the sign of $r$ and whether $r$
is increasing or decreasing determines the nature of the curve's motion
in the plane. \\
\end{tabularx}
\end{skillbox}

\vspace{0.08in}

%----------------------------------------------------------------------
% Vocabulary
%----------------------------------------------------------------------
\begin{skillbox}[{Vocabulary, Concepts, \& Theorems}]{greenbox}
\begin{tabularx}{\linewidth}{@{} p{0.30\linewidth} X @{}}
\textbf{Term} & \textbf{Definition / Note} \\
\midrule
rate of change of $r$ &
  How quickly the radius changes per unit change in angle $\theta$; the
  average rate of change over $[\theta_1, \theta_2]$ is
  $(f(\theta_2)-f(\theta_1))/(\theta_2-\theta_1)$. \\[4pt]
positive and increasing &
  $r > 0$ and $r$ is growing --- the curve moves \emph{away} from the origin
  as $\theta$ increases. \\[4pt]
positive and decreasing &
  $r > 0$ and $r$ is shrinking --- the curve moves \emph{toward} the origin
  as $\theta$ increases. \\[4pt]
negative and increasing &
  $r < 0$ and its absolute value is shrinking --- the curve (on the opposite
  ray) moves toward the origin. \\[4pt]
negative and decreasing &
  $r < 0$ and its absolute value is growing --- the curve (on the opposite
  ray) moves away from the origin. \\
\end{tabularx}
\end{skillbox}

\vspace{0.08in}

%----------------------------------------------------------------------
% Activate Prior Knowledge
%----------------------------------------------------------------------
\begin{skillbox}[Activate Prior Knowledge \& Spiral Review]{sky}
\begin{tabularx}{\linewidth}{@{}X c@{}}
\begin{minipage}[t]{0.96\linewidth}\vspace{0pt}
\textbf{Spiral topics activated during Warm-Up (first 8--10 minutes):}
\begin{itemize}[leftmargin=1.4em, itemsep=2pt]
  \item \textbf{Average rate of change}: Students computed average rates of change for
    polynomial and trigonometric functions in earlier units. The formula
    $\frac{f(b)-f(a)}{b-a}$ carries directly to polar functions.
  \item \textbf{Evaluating trigonometric functions}: Students evaluate $\sin\theta$ and
    $\cos\theta$ at key angles to build tables for polar functions such as
    $r = 2\sin\theta$ or $r = 1 + \cos\theta$.
  \item \textbf{Increasing vs.\ decreasing from a table}: Students identify whether a
    function is increasing or decreasing over an interval by comparing successive
    output values in a table of values.
\end{itemize}
\end{minipage}
&
\end{tabularx}
\end{skillbox}

\vspace{0.08in}

%----------------------------------------------------------------------
% Hook
%----------------------------------------------------------------------
\begin{skillbox}[Hook --- Entry Event]{sky}
\textbf{Scenario:} Display the following table of values for $r = 1 + \sin\theta$
on the board or projector:

\vspace{6pt}
\renewcommand{\arraystretch}{1.4}
\begin{tabular}{|c|c|c|c|c|c|c|c|c|}
\hline
$\theta$ & $0$ & $\pi/4$ & $\pi/2$ & $3\pi/4$ & $\pi$ & $5\pi/4$ & $3\pi/2$ & $7\pi/4$ \\
\hline
$r$ & 1 & $\approx 1.71$ & 2 & $\approx 1.71$ & 1 & $\approx 0.29$ & 0 & $\approx 0.29$ \\
\hline
\end{tabular}

\vspace{8pt}
\textbf{Discussion prompts:}
\begin{enumerate}[leftmargin=1.4em, itemsep=3pt]
  \item \textit{``On which intervals in the table is $r$ getting larger? On which
    is $r$ getting smaller?''}
  \item \textit{``When $r$ is growing, is the point tracing the curve moving toward
    or away from the origin? What about when $r$ is shrinking?''}
  \item \textit{``What happens at $\theta = 3\pi/2$, where $r = 0$? What does that
    mean geometrically?''}
\end{enumerate}

\textbf{Key tension to surface:} A polar curve is not just a static shape --- it is
traced dynamically as $\theta$ sweeps from $0$ to $2\pi$. The rate at which $r$
changes tells us how quickly the curve spirals in or out. Rates of change, the central
tool of Unit 1, appear here in a polar context.

\textbf{Transition:} ``Today we formalize the concept of \emph{rate of change of $r$}
and learn a four-case classification system that captures the full range of possible
curve behavior as $\theta$ increases.''
\end{skillbox}

\vspace{0.08in}

%----------------------------------------------------------------------
% Lesson
%----------------------------------------------------------------------
\begin{skillbox}[Lesson --- Instruction Sequence]{sky}
\begin{multicols}{2}

\textbf{Step 1 --- Average Rate of Change of $r$} (12 min)

Formalize EK 3.15.A.1. Write the formula:
$$\text{ARC} = \frac{f(\theta_2) - f(\theta_1)}{\theta_2 - \theta_1}$$
Interpret: for every additional radian that $\theta$ increases, $r$ changes by
approximately this amount. Work through $r = 2\sin\theta$:
\begin{itemize}[leftmargin=1.3em, itemsep=2pt]
  \item ARC on $[0, \pi/2]$: $\frac{2\sin(\pi/2) - 2\sin 0}{\pi/2 - 0} = \frac{2}{\pi/2} = \frac{4}{\pi} \approx 1.27$
  \item ARC on $[\pi/2, \pi]$: $\frac{0 - 2}{\pi/2} = -\frac{4}{\pi} \approx -1.27$
\end{itemize}
Discuss sign: positive ARC means $r$ is increasing; negative ARC means $r$ is decreasing.

\columnbreak

\textbf{Step 2 --- The Four Cases} (12 min)

Introduce EKs 3.15.A.2--3.15.A.3. Build the four-case table:

\vspace{4pt}
\renewcommand{\arraystretch}{1.4}
\begin{tabular}{|l|l|}
\hline
\textbf{Situation} & \textbf{Curve behavior} \\
\hline
$r > 0$, increasing & away from origin \\
$r > 0$, decreasing & toward origin \\
$r < 0$, increasing & toward origin (opposite ray) \\
$r < 0$, decreasing & away from origin (opposite ray) \\
\hline
\end{tabular}

\vspace{4pt}
Emphasize: when $r < 0$ the point is plotted on the \emph{opposite} terminal ray.
A ``less negative'' $r$ (increasing) means the point is actually getting closer to
the origin on that opposite ray.

\end{multicols}

\begin{multicols}{2}

\textbf{Step 3 --- Reading Intervals from a Table} (10 min)

Use the hook table for $r = 1 + \cos\theta$ at eight key angles. Students
identify by reading the table (no graphing) all intervals where $r$ is
positive and increasing, positive and decreasing, zero, or negative (the
cardioid $r = 1 + \cos\theta$ has $r \geq 0$ for all $\theta$, so focus
on increasing/decreasing and the zero at $\theta = \pi$).

Guided question: \textit{``How can you confirm that $r$ is increasing on $[0, \pi/2]$
without sketching the graph?''} (Compare successive table values.)

\columnbreak

\textbf{Step 4 --- Connecting to Curve Shape} (8 min)

Explain informally how the four-case classification connects to familiar polar
curve shapes. For a cardioid ($r = 1 + \cos\theta$): on $[0, \pi]$ the curve
traces from the rightmost point $(2, 0)$ toward the origin (pole) as $r$ shrinks
from 2 to 0. On $[\pi, 2\pi]$ it traces back outward. For a limaçon with inner
loop ($r = 1 - 2\sin\theta$), $r$ goes negative, tracing the inner loop on the
opposite ray.

\end{multicols}

\smallskip
\textbf{Pacing note:} Steps 1--2 in Period 1; Steps 3--4 and group
activity in Period 2; review and exit ticket in Period 3 if needed.
\end{skillbox}

\vspace{0.08in}

%----------------------------------------------------------------------
% Active Monitoring
%----------------------------------------------------------------------
\begin{skillbox}[Active Monitoring]{redbox}
\begin{itemize}[leftmargin=1.4em, itemsep=3pt]
  \item \textbf{Confusing ARC sign with curve direction (most common)}: Students
    sometimes flip the meaning. Prompt: \textit{``If ARC $> 0$, is $r$ growing or
    shrinking? What does that mean about the distance from the origin?''}
  \item \textbf{Negative $r$ interpretation}: Students may not believe a point can be
    plotted ``backwards.'' Demonstrate with $r = -1$ at $\theta = 0$: plotted at angle
    $\theta + \pi = \pi$, distance 1. The opposite-ray rule is key.
  \item \textbf{ARC units}: Students may omit units or forget the denominator is a
    change in angle (radians). Reinforce: ARC has units of
    \textit{(units of $r$) per radian}.
  \item \textbf{Increasing vs.\ positive}: Students confuse ``$r$ is positive'' with
    ``$r$ is increasing.'' Prompt: \textit{``$r = 0.5$ to $r = 0.2$ is positive but
    \emph{decreasing}. What does that tell you about the curve?''}
\end{itemize}
\end{skillbox}

\vspace{0.08in}

%----------------------------------------------------------------------
% Group Work & Differentiation
%----------------------------------------------------------------------
\begin{skillbox}[Group Work \& Differentiation]{redbox}
Students work in groups of 3--4 on the tiered activity (teacher-assigned
or student self-selected with guidance). All tiers use provided tables ---
no graphing required.

\begin{itemize}[leftmargin=1.4em, itemsep=4pt]
  \item \textbf{Tier R (Reinforce)}: Given a complete table for $r = 2\cos\theta$
    at nine angles, students classify each adjacent interval as one of the four
    cases, then compute the average rate of change over $[0, \pi/2]$ and
    $[\pi/2, \pi]$.
  \item \textbf{Tier A (Apply)}: Given tables for both $r = 2\cos\theta$ and
    $r = 1 - 2\sin\theta$, students identify all intervals where $r < 0$, compute
    average rates of change over two specified intervals for each function, and
    interpret one rate of change in the radar context.
  \item \textbf{Tier E (Extend)}: Given $r = \sin(2\theta)$ (the four-petal rose),
    students use a provided table to classify all intervals in $[0, 2\pi]$ by sign
    and direction, determine which intervals produce each petal, and write an
    explanation of how the sign of $r$ determines which direction each petal is traced.
\end{itemize}
\end{skillbox}

\vspace{0.08in}

%----------------------------------------------------------------------
% Individual Work & Assessment
%----------------------------------------------------------------------
\begin{skillbox}[Individual Work \& Assessment]{redbox}
Exit ticket administered independently (final 5--7 minutes of Period 2 or 3).

\begin{enumerate}[leftmargin=1.4em, itemsep=4pt]
  \item Given a partial table for $r = 1 + \sin\theta$, students classify two
    intervals as positive-increasing or positive-decreasing.
  \item Students compute the average rate of change of $r$ from $\theta = \pi/2$
    to $\theta = \pi$.
  \item Students write one sentence explaining what a negative rate of change of
    $r$ means for a polar function.
\end{enumerate}

\textbf{Data use:} Sort exit tickets by error type --- ARC sign confusion,
negative-$r$ interpretation errors. Use results to determine whether a
brief re-teach opener is needed at the start of Unit 4.
\end{skillbox}

\vspace{0.08in}

%----------------------------------------------------------------------
% Reinforcement & Extension
%----------------------------------------------------------------------
\begin{skillbox}[Reinforcement \& Extension]{goldbox}
\textbf{Homework (Lesson 3.15):} 5--6 problems covering: reading a provided table
for $r = 3\sin\theta$ to classify intervals; computing 2--3 average rates of change;
interpreting a rate of change in context; connecting positive/negative $r$ to curve
direction; 2 AP-style MC items.

\textbf{Preview of Unit 4:} This is the final lesson of Unit 3.
Ask students: \textit{``Rates of change appeared here with polar functions; in
Unit 4 (Functions Involving Parameters, Vectors, and Matrices) we will encounter
parametric functions where both $x$ and $y$ change with a parameter $t$. How might
rates of change apply there?''} Preview the idea that average rate of change of
$x$ or $y$ with respect to $t$ will give direction and speed information for a
parametric curve --- a natural continuation of today's ideas.
\end{skillbox}

\end{document}
